\subsection{Extension of a measure defined on a family of sets}

Let \(\Phi\) be a nonempty set of subsets of a locally compact space \(X\) . Given a real-valued function \(M \mapsto \alpha(M)\) , defined and \(\geqslant 0\) on \(\Phi\) , we propose to seek conditions under which there exists a positive measure \(\mu\) on \(X\) such that the sets of \(\Phi\) are \(\mu\) -integrable and \(\mu(M) = \alpha(M)\) for all \(M \in \Phi\) . We shall limit ourselves to considering the case that the set \(\Phi\) satisfies the following conditions:

\((\mathrm{PC}_{\mathrm{I}})\) The union and intersection of two sets of \(\Phi\) belong to \(\Phi\) .

\((\mathrm{PC}_{\mathrm{II}})\) For every pair consisting of a compact set \(K\) and an open set \(U\) in \(X\) such that \(K \subset U\) , there exists a set \(M \in \Phi\) such that \(K \subset M \subset U\) .

Note that the condition \((\mathrm{PC}_{\mathrm{II}})\) implies that \(\varnothing\in\Phi\) , by taking \(K=U=\varnothing\) . However, the set \(\Phi\) is not necessarily a clan; for example, the set of all compact subsets of X satisfies the conditions \((\mathrm{PC}_{\mathrm{I}})\) and \((\mathrm{PC}_{\mathrm{II}})\) , but in general is not a clan, because if M and N are compact, the same is not in general true of \(M\cap \complement N\) .

We shall assume in addition that the function \(\alpha\) defined on \(\Phi\) satisfies the following conditions (obviously necessary for the problem to have a solution):

\((\mathrm{PM}_{\mathrm{I}})\) The relation \(\mathrm{M}\subset \mathrm{N}\) implies \(\alpha (\mathrm{M})\leqslant \alpha (\mathrm{N})\)

\((\mathrm{PM}_{\mathrm{II}})\) For any M and N in \(\Phi\) , \(\alpha (\mathrm{M}\cup \mathrm{N})\leqslant \alpha (\mathrm{M}) + \alpha (\mathrm{N})\)

\((\mathrm{PM}_{\mathrm{III}})\) The relation \(\mathrm{M}\cap \mathrm{N} = \emptyset\) implies \(\alpha (\mathrm{M}\cup \mathrm{N}) = \alpha (\mathrm{M}) + \alpha (\mathrm{N})\)

On taking \(\mathbf{N} = \varnothing\) in the condition \((\mathrm{PM}_{\mathrm{III}})\) , one deduces that \(\alpha(\varnothing) = 0\) ; the condition \((\mathrm{PM}_{\mathrm{I}})\) then shows that \(\alpha(M) \geqslant 0\) for every \(M \in \Phi\) .

THEOREM 5. --- Let \(\Phi\) be a set of subsets of a locally compact space \(X\) , satisfying \((\mathrm{PC}_{\mathrm{I}})\) and \((\mathrm{PC}_{\mathrm{II}})\) , and let \(\alpha\) be a real-valued function, defined on \(\Phi\) , satisfying the conditions \((\mathrm{PM}_{\mathrm{I}})\) , \((\mathrm{PM}_{\mathrm{II}})\) and \((\mathrm{PM}_{\mathrm{III}})\) . In order that there exist a positive measure \(\mu\) on \(X\) such that the sets of \(\Phi\) are \(\mu\) -integrable and \(\mu(M) = \alpha(M)\) for all \(M \in \Phi\) , it is necessary and sufficient that \(\alpha\) satisfy in addition the following condition:

\((PM_{IV})\) For every \(\varepsilon > 0\) and every \(M \in \Phi\) , there exist a compact set \(K \subset M\) and an open set \(U \supset M\) such that, for every \(N \in \Phi\) satisfying the relation \(K \subset N \subset U\) , one has \(|\alpha(N) - \alpha(M)| \leqslant \varepsilon\) .

Moreover, if the condition \((\mathrm{PM}_{\mathrm{IV}})\) is satisfied, then the measure \(\mu\) is unique; for every compact set K, \(\mu(\mathrm{K}) = \inf_{\mathrm{M} \in \Phi, \mathrm{M} \supset \mathrm{K}} \alpha(\mathrm{M})\) ; for every open set U, \(\mu^{*}(\mathrm{U}) =\)

sup \(\alpha (\mathbf{M})\)

\(\mathbf{M}\in \Phi ,\bar{\mathbf{M}}\subset \mathbf{U}\)

Note that the condition \((PM_{IV})\) is equivalent to the conjunction of the following two:

\((\mathrm{PM}_{\mathrm{IV}}^{\prime})\) For every \(\varepsilon > 0\) and every \(\mathbf{M} \in \Phi\) , there exists an open set \(\mathbf{U} \supset \mathbf{M}\) such that, for every \(\mathbf{N} \in \Phi\) contained in \(\mathbf{U}\) , \(\alpha(\mathbf{N}) \leqslant \alpha(\mathbf{M}) + \varepsilon\) . \((\mathrm{PM}_{\mathrm{IV}}^{\prime \prime})\) For every \(\varepsilon > 0\) and every \(\mathbf{M} \in \Phi\) , there exists a compact set \(\mathbf{K} \subset \mathbf{M}\) such that, for every \(\mathbf{N} \in \Phi\) containing \(\mathbf{K}\) , \(\alpha(\mathbf{N}) \geqslant \alpha(\mathbf{M}) - \varepsilon\) .

For, it is obvious that \((\mathrm{PM}_{\mathrm{IV}}^{\prime})\) and \((\mathrm{PM}_{\mathrm{IV}}^{\prime\prime})\) imply \((\mathrm{PM}_{\mathrm{IV}})\) . Conversely, let us show for example that \((\mathrm{PM}_{\mathrm{IV}})\) implies \((\mathrm{PM}_{\mathrm{IV}}^{\prime})\) : let K be a compact set and U an open set such that \(K \subset M \subset U\) and \(|\alpha(P) - \alpha(M)| \leqslant \varepsilon\) for every \(P \in \Phi\) satisfying \(K \subset P \subset U\) . Then, if \(N \in \Phi\) and \(N \subset U\) , \(M \cup N\) belongs to \(\Phi\) and \(K \subset M \cup N \subset U\) , whence \(\alpha(M \cup N) \leqslant \alpha(M) + \varepsilon\) and a fortiori \(\alpha(N) \leqslant \alpha(M) + \varepsilon\) .

When the set \(\Phi\) , satisfying \((\mathrm{PC}_{1})\) and \((\mathrm{PC}_{11})\) , consists of compact sets, then the condition \((\mathrm{PM}_{\mathrm{IV}}^{\prime \prime})\) is trivially verified, and \((\mathrm{PM}_{\mathrm{IV}})\) is then equivalent to \((\mathrm{PM}_{\mathrm{IV}}^{\prime})\) .

The condition \((\mathrm{PM}_{\mathrm{IV}})\) is necessary: this follows at once from Th. 4 of No.~6 on the `approximation' of an integrable set by a compact set and an open set. To prove the other assertions of the theorem, we proceed in several steps.

\(1^{\circ}\) Definition of a topology on \(\mathfrak{P}(\mathrm{X})\) .

For every pair \((K,U)\) consisting of a compact set K and an open set U in X, we denote by \(I(K,U)\) the set of subsets \(M \subset X\) such that \(K \subset M \subset U\) ; in order that \(I(K,U)\) be nonempty, it is necessary and sufficient that \(K \subset U\) . If \((K',U')\) is a second pair, formed by a compact set \(K'\) and an open set \(U'\) , we have

\[
\mathrm{I} (\mathrm{K}, \mathrm{U}) \cap \mathrm{I} (\mathrm{K} ^ {\prime}, \mathrm{U} ^ {\prime}) = \mathrm{I} (\mathrm{K} \cup \mathrm{K} ^ {\prime}, \mathrm{U} \cap \mathrm{U} ^ {\prime}).
\]

Let \(\mathcal{T}\) be the topology on \(\mathfrak{P}(\mathrm{X})\) generated by the set of subsets I(K,U) as K runs over the set of compact subsets of X, and U over the set of open subsets of X; by the foregoing, the I(K,U) form a base for the topology \(\mathcal{T}\) (GT, I, §1, No.~3).

We observe that the definition of \(\mathcal{T}\) implies that, in \(\mathfrak{P}(\mathrm{X})\) , the set of compact subsets of \(\mathrm{X}\) is dense. The condition \((\mathrm{PC}_{\mathrm{II}})\) expresses that \(\Phi\) is dense in \(\mathfrak{P}(\mathrm{X})\) , and condition \((\mathrm{PM}_{\mathrm{IV}})\) expresses that the function \(\alpha\) is continuous on \(\Phi\) for the topology induced by \(\mathcal{T}\) . Finally, Th. 4 of No.~6 expresses that the function \(\mathrm{M} \mapsto \mu(\mathrm{M})\) is continuous on the clan of \(\mu\) -integrable sets, for the topology induced by \(\mathcal{T}\) .

\(2^{\circ}\) Uniqueness of \(\mu\) .

We denote by \(\overline{\Phi}\) the set of subsets \(M \subset X\) such that \(\alpha(N)\) tends to a finite limit as N tends to M (for the topology T) while remaining in \(\Phi\) ; we may then extend \(\alpha\) in only one way to a continuous mapping \(\overline{\alpha}\) of \(\overline{\Phi}\) into R (GT, I, §8, No.~5, Th. 1). If there exists a measure \(\mu\) meeting the requirements, the above remarks prove that the clan \(\Psi\) of \(\mu\) -integrable sets is contained in \(\overline{\Phi}\) and that \(\mu(M) = \overline{\alpha}(M)\) for every \(M \in \Psi\) ; this relation holds in particular for every compact subset M of X, which proves the uniqueness of \(\mu\) (No.~10, Cor. 3 of Prop. 19).

\(3^{\circ}\) Extension of \(\alpha\) to the compact sets.

Without assuming the existence of \(\mu\) , we are now going to study the set \(\overline{\Phi}\) and the extension \(\overline{\alpha}\) of \(\alpha\) to \(\overline{\Phi}\) . We first show that every compact set \(K\) belongs to \(\overline{\Phi}\) and that \(\overline{\alpha}(K) = \inf_{P \in \Phi, P \supset K} \alpha(P)\) . Set \(a = \inf_{P \in \Phi, P \supset K} \alpha(P)\) ; for every \(\varepsilon > 0\) , there exists an \(M \in \Phi\) such that \(K \subset M\) and \(\alpha(M) \leqslant a + \varepsilon\) . By \((PM_{IV}')\) , there exists an open set \(U \supset M\) such that, for every \(N \in \Phi\) contained in \(U\) , we have \(\alpha(N) \leqslant \alpha(M) + \varepsilon \leqslant a + 2\varepsilon\) ; for every \(N \in \Phi\) such that \(K \subset N \subset U\) , we therefore have \(a \leqslant \alpha(N) \leqslant a + 2\varepsilon\) , which, by the definitions, shows that \(K \in \overline{\Phi}\) and \(\overline{\alpha}(K) = a\) .

This result proves at once that if \(K_{1}\) and \(K_{2}\) are two compact sets such that \(K_{1} \subset K_{2}\) , then \(\overline{\alpha}(K_{1}) \leqslant \overline{\alpha}(K_{2})\) . If \(K_{1}\) and \(K_{2}\) are any two compact sets, we have \(\overline{\alpha}(K_{1} \cup K_{2}) \leqslant \overline{\alpha}(K_{1}) + \overline{\alpha}(K_{2})\) by \((\mathrm{PM}_{\mathrm{II}})\) . We shall see, moreover, that if \(K_{1}\) and \(K_{2}\) are disjoint then \(\overline{\alpha}(K_{1} \cup K_{2}) = \overline{\alpha}(K_{1}) + \overline{\alpha}(K_{2})\) . For, there then exist two disjoint open sets \(U_{1}, U_{2}\) such that \(K_{1} \subset U_{1}, K_{2} \subset U_{2}\) (GT, II, §4, Prop. 4). Therefore, by \((\mathrm{PC}_{\mathrm{II}})\) , there also exist two sets \(M_{1} \in \Phi, M_{2} \in \Phi\) such that \(K_{1} \subset M_{1} \subset U_{1}\) and \(K_{2} \subset M_{2} \subset U_{2}\) . Now let P be any set of \(\Phi\) containing \(K_{1} \cup K_{2}\) ; the union of the two sets \(P \cap M_{1}\) and \(P \cap M_{2}\) belongs to \(\Phi\) by \((\mathrm{PC}_{\mathrm{I}})\) , and since these two sets are disjoint, application of \((\mathrm{PM}_{\mathrm{I}})\) and \((\mathrm{PM}_{\mathrm{III}})\) yields

\[
\alpha (\mathrm{P}) \geqslant \alpha (\mathrm{P} \cap \mathrm{M} _ {1}) + \alpha (\mathrm{P} \cap \mathrm{M} _ {2}) \geqslant \overline {{\alpha}} (\mathrm{K} _ {1}) + \overline {{\alpha}} (\mathrm{K} _ {2}),
\]

which establishes our assertion.

\(4^{\circ}\) Extension of \(\alpha\) to the open sets.

We shall now see that, for an open set U to belong to \(\overline{\Phi}\) , it is necessary and sufficient that, as K runs over the set of compact subsets of U, the supremum of the numbers \(\overline{\alpha}(K)\) is finite; moreover, \(\overline{\alpha}(U)\) is then equal to this supremum.

For, let U be an open set belonging to \(\bar{\Phi}\) ; for every \(\varepsilon > 0\) there exists a compact set \(K \subset U\) such that, for every set \(M \in \Phi\) satisfying \(K \subset M \subset U\) , one has \(|\overline{\alpha}(U) - \alpha(M)| \leqslant \varepsilon\) , whence \(|\overline{\alpha}(U) - \overline{\alpha}(K)| \leqslant \varepsilon\) ; on the other hand, if \(K'\) is any compact set contained in U, then \(K \subset K \cup K' \subset U\) , whence \(|\overline{\alpha}(U) - \overline{\alpha}(K \cup K')| \leqslant \varepsilon\) and so \(\overline{\alpha}(U) \geqslant \overline{\alpha}(K \cup K') - \varepsilon \geqslant \overline{\alpha}(K') - \varepsilon\) ; \(\overline{\alpha}(U)\) is therefore indeed equal to the supremum of the numbers \(\overline{\alpha}(K)\) as K runs over the set of compact subsets of U.

Conversely, let U be an open set such that \(b = \sup_{K \subset U} \overline{\alpha}(K) < +\infty\) (K run-

ning over the set of compact subsets of U), and let us show that \(U \in \overline{\Phi}\) . For every \(\varepsilon > 0\) , there exists a compact set \(K \subset U\) such that \(b - \varepsilon \leqslant \overline{\alpha}(K) \leqslant b\) ; by \((\mathrm{PM}_{\mathrm{IV}}^{\prime\prime})\) , for every set \(M \in \Phi\) such that \(K \subset M \subset U\) , there exists a compact set \(K' \subset M\) such that

\[
\alpha (\mathrm{M}) \leqslant \overline {{{{\alpha}}}} (\mathrm{K} ^ {\prime}) + \varepsilon \leqslant b + \varepsilon ;
\]

therefore \(b - \varepsilon \leqslant \alpha (\mathrm{M}) \leqslant b + \varepsilon\) , which proves that \(\mathrm{U} \in \overline{\Phi}\) .

From this characterization of the open sets \(U \in \overline{\Phi}\) , and of \(\overline{\alpha}(U)\) , it follows first of all that if \(U_{1}\) and \(U_{2}\) are two open sets such that \(U_{1} \subset U_{2}\) and \(U_{2} \in \overline{\Phi}\) , then \(U_{1} \in \overline{\Phi}\) and \(\overline{\alpha}(U_{1}) \leqslant \overline{\alpha}(U_{2})\) . On the other hand, if \(U_{1}\) and \(U_{2}\) are two open sets belonging to \(\overline{\Phi}\) , then the same is true of \(U_{1} \cup U_{2}\) , and \(\overline{\alpha}(U_{1} \cup U_{2}) \leqslant \overline{\alpha}(U_{1}) + \overline{\alpha}(U_{2})\) . For, let K be any compact set contained in \(U_{1} \cup U_{2}\) ; for every point \(x \in K\) , there exists a compact neighborhood of x contained in either \(U_{1}\) or \(U_{2}\) ; one can therefore cover K by a finite number of these neighborhoods; if \(K_{1}\) (resp. \(K_{2}\) ) is the union of those that are contained in \(U_{1}\) (resp. \(U_{2}\) ), then \(K \subset K_{1} \cup K_{2}\) , whence

\[
\overline {{{{\alpha}}}} (\mathrm{K}) \leqslant \overline {{{{\alpha}}}} (\mathrm{K} _ {1} \cup \mathrm{K} _ {2}) \leqslant \overline {{{{\alpha}}}} (\mathrm{K} _ {1}) + \overline {{{{\alpha}}}} (\mathrm{K} _ {2}) \leqslant \overline {{{{\alpha}}}} (\mathrm{U} _ {1}) + \overline {{{{\alpha}}}} (\mathrm{U} _ {2}),
\]

which establishes the asserted property.

\(5^{\circ}\) Properties of \(\overline{\Phi}\) and \(\overline{\alpha}\) .

The definition of \(\overline{\Phi}\) and \(\overline{\alpha}\) can now be transformed as follows (taking into account \((\mathrm{PC}_{\mathrm{II}})\) ): in order that \(M \in \overline{\Phi}\) , it is necessary and sufficient that, for every \(\varepsilon > 0\) , there exist a compact set K and an open set \(U \in \overline{\Phi}\) such that \(K \subset M \subset U\) and \(\overline{\alpha}(U) - \overline{\alpha}(K) \leqslant \varepsilon\) ; \(\overline{\alpha}(M)\) is, moreover, the infimum of the \(\overline{\alpha}(U)\) for the open sets \(U \in \overline{\Phi}\) containing M, and the supremum of the \(\overline{\alpha}(K)\) for the compact sets \(K \subset M\) .

From this, we shall first deduce that if \(\mathrm{M}_1\) , \(\mathrm{M}_2\) and \(\mathrm{M}_1 \cup \mathrm{M}_2\) belong to \(\overline{\Phi}\) , then \(\overline{\alpha}(\mathrm{M}_1 \cup \mathrm{M}_2) \leqslant \overline{\alpha}(\mathrm{M}_1) + \overline{\alpha}(\mathrm{M}_2)\) . Indeed, if \(\mathrm{U}_1\) and \(\mathrm{U}_2\) are two open sets of \(\overline{\Phi}\) containing \(\mathbf{M}_1\) and \(\mathbf{M}_2\) , respectively, and such that \(\overline{\alpha} (\mathrm{U}_1)\leqslant \overline{\alpha} (\mathrm{M}_1) + \varepsilon\) and \(\overline{\alpha} (\mathrm{U}_2)\leqslant \overline{\alpha} (\mathrm{M}_2) + \varepsilon\) , then \(\mathrm{U}_1\cup \mathrm{U}_2\) belongs to \(\overline{\Phi}\) , contains \(\mathbf{M}_1\cup \mathbf{M}_2\) , and consequently

\[
\overline {{\alpha}} \left(\mathrm{M} _ {1} \cup \mathrm{M} _ {2}\right) \leqslant \overline {{\alpha}} \left(\mathrm{U} _ {1} \cup \mathrm{U} _ {2}\right) \leqslant \overline {{\alpha}} \left(\mathrm{U} _ {1}\right) + \overline {{\alpha}} \left(\mathrm{U} _ {2}\right) \leqslant \overline {{\alpha}} \left(\mathrm{M} _ {1}\right) + \overline {{\alpha}} \left(\mathrm{M} _ {2}\right) + 2 \varepsilon ,
\]

whence our assertion.

Next, let us show that if \(K\) is a compact set and \(U\) is an open set of \(\overline{\Phi}\) such that \(K \subset U\) , then \(\overline{\alpha}(U - K) = \overline{\alpha}(U) - \overline{\alpha}(K)\) . By the foregoing, we have \(\overline{\alpha}(U) \leqslant \overline{\alpha}(K) + \overline{\alpha}(U - K)\) . On the other hand, for every compact set \(K' \subset U - K\) ,

\[
\overline {{{{\alpha}}}} (\mathrm{K} \cup \mathrm{K} ^ {\prime}) = \overline {{{{\alpha}}}} (\mathrm{K}) + \overline {{{{\alpha}}}} (\mathrm{K} ^ {\prime}) \leqslant \overline {{{{\alpha}}}} (\mathrm{U});
\]

since U - K is open and belongs to \(\overline{\Phi}\) , \(\overline{\alpha}(U - K)\) is the supremum of the \(\overline{\alpha}(K')\) , which shows that \(\overline{\alpha}(K) + \overline{\alpha}(U - K) \leqslant \overline{\alpha}(U)\) .

The definition of \(\overline{\Phi}\) may therefore now be expressed in the following way: in order that \(M \in \overline{\Phi}\) , it is necessary and sufficient that, for every \(\varepsilon > 0\) , there exist a compact set K and an open set \(U \in \overline{\Phi}\) such that \(K \subset M \subset U\) and \(\overline{\alpha}(U - K) \leqslant \varepsilon\) .

We are now in a position to prove that \(\overline{\Phi}\) is a clan and \(\overline{\alpha}\) an additive set function on \(\overline{\Phi}\) . We first show that if M and N belong to \(\overline{\Phi}\) then so do M \(\cap\) CN and M \(\cup\) N. By hypothesis, for every \(\varepsilon > 0\) there exist two compact sets K,K' and two open sets U,U' of \(\overline{\Phi}\) such that

\[
\mathrm{K} \subset \mathrm{M} \subset \mathrm{U}, \quad \mathrm{K} ^ {\prime} \subset \mathrm{N} \subset \mathrm{U} ^ {\prime}, \quad \overline {{\alpha}} (\mathrm{U} - \mathrm{K}) \leqslant \varepsilon , \quad \overline {{\alpha}} (\mathrm{U} ^ {\prime} - \mathrm{K} ^ {\prime}) \leqslant \varepsilon .
\]

The set \(K'' = K \cap C U'\) is compact, the set \(U'' = U \cap C K'\) is open and belongs to \(\overline{\Phi}\) , and \(K'' \subset M \cap C N \subset U''\) ; on the other hand, \(U'' - K''\) is contained in the union of \(U \cap C K\) and \(U' \cap C K'\) , whence \(\overline{\alpha}(U'' - K'') \leqslant 2\varepsilon\) , which proves that \(M \cap C N \in \overline{\Phi}\) . Similarly, \(U_1 = U \cup U'\) is open and belongs to \(\overline{\Phi}\) , \(K_1 = K \cup K'\) is compact, and \(K_1 \subset M \cup N \subset U_1\) ; on the other hand, \(U_1 - K_1\) is contained in the union of U - K and \(U' - K'\) , whence again \(\overline{\alpha}(U_1 - K_1) \leqslant 2\varepsilon\) , and \(M \cup N\) belongs to \(\overline{\Phi}\) . Finally, if M and N are disjoint, then

\[
\overline {{{{\alpha}}}} \left(\mathrm{K} _ {1}\right) = \overline {{{{\alpha}}}} (\mathrm{K}) + \overline {{{{\alpha}}}} \left(\mathrm{K} ^ {\prime}\right) \geqslant \overline {{{{\alpha}}}} (\mathrm{M}) + \overline {{{{\alpha}}}} (\mathrm{N}) - 2 \varepsilon ,
\]

consequently \(\overline{\alpha} (\mathbf{M}\cup \mathbf{N})\geqslant \overline{\alpha} (\mathbf{M}) + \overline{\alpha} (\mathbf{N}) - 2\varepsilon\) since \(\varepsilon\) is arbitrary, we have \(\overline{\alpha} (\mathbf{M}\cup \mathbf{N}) = \overline{\alpha} (\mathbf{M}) + \overline{\alpha} (\mathbf{N})\)

\(6^{\circ}\) Existence of the measure \(\mu\) .

By Prop. 18 of No.~9, there exists one and only one positive linear form \(\beta\) on the vector space \(\mathcal{E}(\overline{\Phi})\) of \(\overline{\Phi}\) -step functions, such that \(\beta(\varphi_{\mathrm{M}}) = \overline{\alpha}(\mathrm{M})\) for all \(\mathrm{M} \in \overline{\Phi}\) . For every compact subset \(\mathrm{K}\) of \(\mathrm{X}\) , let us denote by \(\mathcal{G}(\mathrm{K})\) the space of uniform limits of functions of \(\mathcal{E}(\overline{\Phi})\) whose support is contained in \(\mathrm{K}\) . Since \(\beta\) is positive, \(|\beta(f)| \leqslant \overline{\alpha}(\mathrm{K}) \cdot \|f\|\) for every function \(f \in \mathcal{E}(\overline{\Phi})\) whose support is contained in \(\mathrm{K}\) ; the restriction of \(\beta\) to the space of these functions is a continuous linear form for the topology of uniform convergence; it may therefore be extended to a positive continuous linear form \(\overline{\beta}_{K}\) on \(\mathcal{G}(K)\) . Moreover, if K and \(K_{1}\) are two compact sets such that \(K \subset K_{1}\) , then the restriction of \(\overline{\beta}_{K_{1}}\) to \(\mathcal{G}(K)\) is identical to \(\overline{\beta}_{K}\) , therefore there exists a positive linear form \(\overline{\beta}\) on the union G of the \(\mathcal{G}(K)\) , that extends each of the forms \(\overline{\beta}_{K}\) .

Now, since every compact set belongs to \(\overline{\Phi}\) , the space \(\mathcal{K}\) of continuous real-valued functions with compact support is a subspace of \(\mathcal{G}\) (No.~10, Prop. 19); the restriction to \(\mathcal{K}\) of the positive linear form \(\overline{\beta}\) is therefore a positive measure \(\mu\) . Let us show that for every compact set \(K\) , \(\mu(K) = \overline{\alpha}(K)\) . For every \(\varepsilon > 0\) , there exists an open set \(U \in \overline{\Phi}\) such that \(K \subset U\) , \(\mu(U) \leqslant \mu(K) + \varepsilon\) and \(\overline{\alpha}(U) \leqslant \overline{\alpha}(K) + \varepsilon\) . Let \(f\) be a continuous mapping of \(X\) into {[}0,1{]} whose support is contained in \(U\) and such that \(f(x) = 1\) on \(K\) (Ch. III, §1, No.~2, Lemma 1). Then \(\mu(K) \leqslant \mu(f) \leqslant \mu(U) \leqslant \mu(K) + \varepsilon\) , and, on the other hand,

\[
\overline {{\alpha}} (\mathrm{K}) = \beta (\varphi_ {\mathrm{K}}) \leqslant \overline {{\beta}} (f) \leqslant \beta (\varphi_ {\mathrm{U}}) = \overline {{\alpha}} (\mathrm{U}) \leqslant \overline {{\alpha}} (\mathrm{K}) + \varepsilon ;
\]

since \(\mu(f) = \overline{\beta}(f)\) , we see that \(|\mu(\mathrm{K}) - \overline{\alpha}(\mathrm{K})| \leqslant \varepsilon\) , and since \(\varepsilon\) is arbitrary, \(\mu(\mathrm{K}) = \overline{\alpha}(\mathrm{K})\) .

The characterization of the open sets belonging to \(\overline{\Phi}\) , combined with Cor. 4 of Th. 4 of No.~6, then shows that the open sets belonging to \(\overline{\Phi}\) are none other than the \(\mu\) -integrable open sets, and that, for such a set U, we have \(\mu(\mathrm{U}) = \overline{\alpha}(\mathrm{U})\) . Th. 4 of No.~6 and the characterization of the sets of \(\overline{\Phi}\) given in \(5^{\circ}\) then show that the \(\mu\) -integrable sets are the sets of \(\overline{\Phi}\) and that, for such a set M, \(\mu(\mathrm{M}) = \overline{\alpha}(\mathrm{M})\) . Finally, the fact that \(\mu^{*}(\mathrm{U}) = \sup_{\mathrm{M} \in \Phi, \mathrm{M} \subset \mathrm{U}} \alpha(\mathrm{M})\) for every open set U follows at once from \((\mathrm{PC}_{\mathrm{II}})\) and Cor. 4 of Th. 4 of No.~6.

Theorem 5 is thus completely proved.

COROLLARY. --- Let X be a locally compact space with a countable base, \(\Psi\) the set of Borel sets of X, \(\beta\) a mapping of \(\Psi\) into \([0,+\infty]\) satisfying the following conditions:

\begin{enumerate}
\def\labelenumi{(\roman{enumi})}
\item
  If \((\mathrm{B}_1, \mathrm{B}_2, \ldots)\) is a sequence of pairwise disjoint Borel sets of X, then \(\beta(\mathrm{B}_1 \cup \mathrm{B}_2 \cup \ldots) = \beta(\mathrm{B}_1) + \beta(\mathrm{B}_2) + \ldots\) .
\item
  If B is a compact subset of X, then \(\beta (\mathrm{B}) < +\infty\)
\end{enumerate}

Then, there exists one and only one positive measure \(\mu\) on \(X\) such that \(\beta(B) = \mu^{*}(B)\) for all \(B \in \Psi\) .

Let \(\Phi\) be the set of compact subsets of X and let \(\alpha\) be the restriction of \(\beta\) to \(\Phi\) . The conditions \((\mathrm{PC_I})\) , \((\mathrm{PC_{II}})\) , \((\mathrm{PM_I})\) , \((\mathrm{PM_{II}})\) , \((\mathrm{PM_{III}})\) and \((\mathrm{PM_{IV}''})\) are then satisfied. Let K be a compact subset of X, and \(\varepsilon > 0\) . Then K is the intersection of a decreasing sequence \((\mathrm{U}_1, \mathrm{U}_2, \ldots)\) of relatively compact open sets of X (GT, IX, §2, No.~5, Prop. 7). We have \(\sum_{n=1}^{\infty} \beta(\mathrm{U}_n - \mathrm{U}_{n+1}) = \beta(\mathrm{U}_1 - \mathrm{K}) < +\infty\) , therefore

\[
\beta (\mathrm{U} _ {n}) - \beta (\mathrm{K}) = \beta (\mathrm{U} _ {n} - \mathrm{K}) = \sum_ {p = n} ^ {\infty} \beta (\mathrm{U} _ {p} - \mathrm{U} _ {p + 1})
\]

tends to 0 as n tends to \(\infty\) . This proves that the condition \((\mathrm{PM}_{\mathrm{IV}}^{\prime})\) is satisfied. By Th. 5, there exists a positive measure \(\mu\) on X such that \(\mu(\mathrm{K}) = \alpha(\mathrm{K})\) for every compact subset K of X. Since every open set U of X is the union of an increasing sequence of compact subsets, we have \(\mu^{*}(\mathrm{U}) = \beta(\mathrm{U})\) . Let L be a compact subset of X. By Prop. 7 of No.~5, the \(\mu\) -integrable subsets of L form a tribe of subsets of L. Therefore, if B is an element of \(\Psi\) contained in L, then B is \(\mu\) -integrable; for every \(\varepsilon > 0\) , there then exist a compact set K and an open set U in X such that \(K \subset B \subset U\) and \(\mu^{*}(\mathrm{U}) - \mu(\mathrm{K}) \leqslant \varepsilon\) (No.~6, Th. 4). Since \(\beta(\mathrm{U}) = \mu^{*}(\mathrm{U})\) and \(\beta(\mathrm{K}) = \mu(\mathrm{K})\) , we see that \(|\mu^{*}(\mathrm{B}) - \beta(\mathrm{B})| \leqslant 2\varepsilon\) . Therefore \(\beta(\mathrm{B}) = \mu^{*}(\mathrm{B})\) . Finally, every Borel set C of X is the union of a sequence of pairwise disjoint, relatively compact Borel sets, whence \(\beta(\mathrm{C}) = \mu^{*}(\mathrm{C})\) . The uniqueness of \(\mu\) follows at once from Th. 5.

